\documentclass[11pt,twocolumn]{article}

\usepackage[margin=0.85in]{geometry}
\usepackage[T1]{fontenc}
\usepackage{amsmath,amssymb}
\usepackage{booktabs}
\usepackage{array}
\usepackage{hyperref}
\usepackage{xcolor}
\usepackage{microtype}
\usepackage{authblk}
\usepackage[numbers]{natbib}
\usepackage{enumitem}
\setlist{nosep}

\hypersetup{colorlinks=true,linkcolor=blue!60!black,citecolor=blue!60!black,urlcolor=blue!60!black}

\title{Synthetic Falsifiers for Recognition Science Predictions\\in Genome Architecture}

\author[1]{Jonathan Washburn}
\affil[1]{Independent Researcher}

\date{March 2026}

\begin{document}
\maketitle

\begin{abstract}
We define a formal framework for testing the Recognition Science (RS)
theory of genome architecture against synthetic biology constructs. Three
machine-checked assay predicates are formalized in Lean~4: (1)~a
\textbf{degeneracy assay} that tests whether a synthetic code preserves the
standard degeneracy pattern; (2)~a \textbf{wobble assay} that tests whether
the clean-wobble family-box count is maintained; and (3)~an
\textbf{optimality assay} that tests whether the synthetic code achieves
lower total weighted mutation cost than the standard code. A synthetic
construct that fails any assay constitutes a formal falsifier of the RS
framework. The standard code is proved to pass all three assays. We
describe five classes of synthetic constructs---reduced-redundancy codes,
phase-misaligned introns, sub-threshold introns, repeat-depleted genomes,
and clock-disrupted timing elements---each designed to probe a specific
layer of the RS five-layer genome decomposition. This is a design paper:
it specifies the constructs, assays, and predicted outcomes, but does not
report experimental results. The purpose is to convert the RS framework
from a post-hoc interpretation into an experimentally falsifiable program.
\end{abstract}

\section{Introduction}

A theoretical framework that cannot be falsified is not science. The RS
theory of genome architecture~\cite{Washburn2025, Paper4} makes specific
structural predictions: the standard genetic code minimizes a
position-weighted mutation cost~\cite{Paper2}, introns are positioned by
helical phase stress~\cite{Paper7}, wobble degeneracy dominates the codon
table~\cite{Paper3}, and the FOLD operator confers a categorical advantage
to sexual reproduction~\cite{Paper6}.

Each of these predictions can in principle be tested by constructing
synthetic codes or genomic architectures that violate the predicted
constraints and measuring the phenotypic consequences. This paper defines
the formal assay framework and specifies five classes of synthetic
falsifier constructs.

The key contribution is that the assay predicates are \textbf{formalized
in Lean~4}. The wobble assay is decidable (equality of natural numbers);
the degeneracy and optimality assays are logical propositions over the
formalized code. The standard code is proved to pass all three assays, so
any construct that fails at least one is a formal falsifier---a
machine-checkable witness that the RS prediction is violated.

\section{The Assay Framework}

\subsection{Synthetic Constructs}

A \textbf{synthetic construct} is a codon assignment---a bijection on the
20 amino acid labels that relabels which amino acid each codon family
encodes. The standard code is the identity assignment:

\begin{quote}
\texttt{standardConstruct} has \texttt{assignment := standardAssignment}
\end{quote}

\noindent (equivalently, \texttt{ofAssignment standardAssignment}).
Both are defined in \texttt{SyntheticFalsifierAssays.lean}.

\subsection{Three Assay Predicates}

\textbf{Assay 1: Degeneracy.}
A construct passes the degeneracy assay if and only if its assignment
respects the standard degeneracy classes:
\[
  \texttt{degeneracyAssay}(s) \;\Leftrightarrow\;
  \texttt{respectsDegeneracyClass}(s.\texttt{assignment}).
\]
A construct that reassigns a 6-fold amino acid to a 2-fold codon family
fails this assay.

\textbf{Assay 2: Wobble.}
A construct passes the wobble assay if the number of family boxes with
clean wobble symmetry (all four third-position variants synonymous, or a
clean pyrimidine/purine $2{+}2$ split) equals the standard count:
\[
  \texttt{wobbleAssay}(s) \;\Leftrightarrow\;
  \texttt{assignedCleanWobbleCount}(s.\texttt{assignment}) =
  \texttt{cleanWobbleCount}.
\]
The standard code has 14 of 16 clean family boxes~\cite{Paper3}. A
reassignment that breaks this count fails the wobble assay.

\textbf{Assay 3: Optimality.}
A construct passes the optimality assay if the standard code's total
weighted mutation cost is no greater than the construct's cost:
\[
  \texttt{optimalityAssay}(w, s) \;\Leftrightarrow\;
  \texttt{standardAssignmentCost}(w) \leq
  \texttt{totalAssignmentCost}(w, s.\texttt{assignment}).
\]
A construct that achieves strictly lower cost than the standard code under
the chosen weight vector $w$ fails the optimality assay.

\subsection{Falsifier Definition}

A construct is \textbf{RS-violating} if it fails at least one of the three
assays:
\[
  \texttt{rsViolatingConstruct}(w, s) \;\Leftrightarrow\;
  \neg\,\texttt{degeneracyAssay}(s) \;\lor\;
  \neg\,\texttt{wobbleAssay}(s) \;\lor\;
  \neg\,\texttt{optimalityAssay}(w, s).
\]
A \texttt{SyntheticFalsifier} packages the construct together with a
proof that it violates at least one assay.

\subsection{The Standard Code Passes All Assays}

\begin{quote}
\textbf{Theorem} (\texttt{standardConstruct\_not\_falsifier}).
For any weight vector $w$, the standard code is not an RS-violating
construct: $\neg\,\texttt{rsViolatingConstruct}(w,
\texttt{standardConstruct})$.
\end{quote}

The proof proceeds by showing that the standard code passes each assay
individually:
\begin{itemize}
\item \texttt{standardConstruct\_passes\_degeneracyAssay}:
  the identity assignment respects all degeneracy classes.
\item \texttt{standardConstruct\_passes\_wobbleAssay}:
  the identity assignment preserves the clean-wobble count.
\item \texttt{standardConstruct\_passes\_optimalityAssay}:
  the identity assignment's cost equals its own cost ($\leq$ is
  reflexive).
\end{itemize}

\section{Five Classes of Falsifier Constructs}

Each class targets a specific layer of the RS five-layer genome
decomposition~\cite{Paper4}.

\subsection{Class 1: Reduced-Redundancy Codes}

\textbf{Target layer}: Error correction.

\textbf{Construct}: Reassign amino acids to reduce average degeneracy
below the standard $61/20 \approx 3.05$. For example, split current
4-fold families into pairs of 2-fold families by introducing new amino
acid assignments at the wobble position.

\textbf{Predicted assay outcome}: Fails the wobble assay (clean-wobble
count drops below 14) and may fail the optimality assay (reduced
redundancy increases vulnerability to point mutations).

\textbf{Falsifier for RS}: If a reduced-redundancy code shows
\emph{improved} fitness in a controlled expression system, the RS
prediction that the current redundancy level is error-optimal is wrong.

\subsection{Class 2: Phase-Misaligned Introns}

\textbf{Target layer}: Phase alignment.

\textbf{Construct}: Insert synthetic introns at positions that violate
the helical phase-stress model: not on the half-octave grid
($p \not\equiv 0 \pmod{4}$), and at low-stress rather than high-stress
sites.

\textbf{Predicted assay outcome}: Reduced splicing efficiency or gene
expression, measurable by RT-qPCR or ribosome profiling.

\textbf{Falsifier for RS}: If phase-misaligned introns produce
\emph{equal or better} expression than phase-aligned introns of the same
length, the helical-phase model~\cite{Paper7} is wrong.

\subsection{Class 3: Sub-Threshold Introns}

\textbf{Target layer}: Phase alignment (minimal intron size).

\textbf{Construct}: Engineer introns shorter than the predicted minimal
size of $8\varphi^5 \approx 89$~bp. Test introns of 50, 60, 70, and
80~bp.

\textbf{Predicted assay outcome}: Introns below $\sim$85~bp should show
impaired splicing or reduced expression in mammalian cells.

\textbf{Falsifier for RS}: If introns of 50--60~bp splice efficiently and
support normal expression in mammalian cells, the $8\varphi^5$
prediction~\cite{Paper7} is wrong.

\subsection{Class 4: Repeat-Depleted Genomes}

\textbf{Target layer}: Ledger balance.

\textbf{Construct}: Systematically delete repetitive elements (SINEs,
LINEs, LTR retrotransposons) from a defined genomic region while
preserving all annotated genes and regulatory elements.

\textbf{Predicted assay outcome}: Impaired gene regulation, chromatin
organization, or replication timing in the depleted region.

\textbf{Falsifier for RS}: If repeat depletion has no measurable
phenotypic consequence (beyond the directly deleted elements), the
ledger-balance hypothesis~\cite{Paper4} is wrong.

\subsection{Class 5: Clock-Disrupted Timing Elements}

\textbf{Target layer}: Temporal coordination.

\textbf{Construct}: Mutate or delete non-coding elements predicted to
implement $\varphi$-ladder timing (e.g., clock gene regulatory regions,
replication-timing control elements) and measure the effect on downstream
timing phenotypes.

\textbf{Predicted assay outcome}: Disrupted biological timing at the
specific $\varphi$-ladder rung corresponding to the deleted element.

\textbf{Falsifier for RS}: If deletion of a predicted clock element has
no effect on any timing phenotype, the gearbox model~\cite{Paper8} is
wrong for that rung.

\section{Formal Verification}

The assay predicates and all structural theorems are formalized in
Lean~4 (version 4.27.0) with Mathlib. The principal module is
\texttt{SyntheticFalsifierAssays.lean}, which imports:

\begin{itemize}
\item \texttt{CodonAssignmentGlobalOptimality.lean}: the
  \texttt{GlobalAssignmentCounterexampleCertificate} type.
\item \texttt{CodonAssignmentOptimality.lean}: the \texttt{swapImproves}
  predicate.
\item \texttt{FamilyBoxAnalysis.lean}: \texttt{allFamilyBoxes},
  \texttt{cleanWobbleCount}, \texttt{hasCleanWobble}.
\item \texttt{DegeneracySearchSpace.lean}:
  \texttt{respectsDegeneracyClass},
  \texttt{standardAssignment\_respectsDegeneracyClass}.
\end{itemize}

Key theorems:
\begin{itemize}
\item \texttt{standardConstruct\_not\_falsifier}: the standard code
  passes all three assays.
\item \texttt{degeneracyFailure\_yields\_falsifier}: any assignment that
  violates degeneracy classes is a falsifier.
\item \texttt{wobbleFailure\_yields\_falsifier}: any assignment that
  changes the clean-wobble count is a falsifier.
\item \texttt{costImprovement\_yields\_falsifier}: any assignment that
  strictly reduces total cost is a falsifier.
\item \texttt{improvingSwap\_yields\_falsifier}: any improving pairwise
  swap is a falsifier.
\item \texttt{globalCounterexample\_yields\_falsifier}: any global
  counterexample certificate is a falsifier.
\end{itemize}

The module contains zero axioms and zero \texttt{sorry} assertions.

\section{Discussion}

\subsection{Why a Design Paper}

This paper specifies constructs, assays, and predictions but does not
report experimental results. The purpose is to convert the RS framework
from a post-hoc interpretation into a pre-registered experimental program.
Each construct class can be independently tested by a synthetic biology
or genomics laboratory without requiring acceptance of the broader RS
theory.

\subsection{Prioritization}

The five construct classes differ in experimental difficulty:

\begin{center}
\small
\begin{tabular}{@{}lll@{}}
\toprule
Class & Difficulty & Technology \\
\midrule
1. Reduced redundancy & Medium & Expanded genetic codes \\
2. Phase-misaligned introns & Low & Synthetic gene design \\
3. Sub-threshold introns & Low & Minigene reporters \\
4. Repeat depletion & High & CRISPR-based deletion \\
5. Clock disruption & High & Enhancer/promoter KO \\
\bottomrule
\end{tabular}
\end{center}

Classes 2 and 3 are the easiest to test and should be prioritized. A
single well-controlled experiment showing that phase-misaligned introns
splice as well as phase-aligned ones would refute the helical-phase model.

\subsection{Relation to Existing Synthetic Code Work}

The expanded genetic code community~\cite{Chin2017} has already
engineered organisms with reassigned codons (e.g., \emph{E.~coli} with
all TAG codons reassigned). These experiments test codon \emph{identity},
not codon \emph{optimality}. The RS falsifier framework asks a different
question: not whether a reassigned code \emph{can work}, but whether it
\emph{outperforms} the standard code on the RS-predicted objective.

\subsection{Scope}

The assay framework currently covers codon-level properties (degeneracy,
wobble, cost). It does not yet cover intron-level or timing-level
properties, which would require extending the Lean formalization with
sequence-level predicates. Classes 2--5 are specified as construct
descriptions, not as formalized assay predicates.

\section{Conclusion}

The RS theory of genome architecture generates specific, testable
predictions at every layer of the five-layer decomposition. This paper
defines a formal assay framework (three machine-checked predicates) and
five classes of synthetic falsifier constructs designed to probe each
layer independently. The standard code is proved to pass all assays. Any
construct that fails at least one assay is a formal, machine-checkable
falsifier. The immediate experimental priorities are phase-misaligned
introns and sub-threshold introns, both of which can be tested with
standard synthetic gene design.

\section*{Data Availability}

All Lean source files are in the \texttt{reality} repository under
\texttt{IndisputableMonolith/Genetics/}.

\begin{thebibliography}{99}
\bibitem{Chin2017}
Chin, J.W. (2017). Expanding and reprogramming the genetic code.
\textit{Nature} \textbf{550}(7674), 53--60.

\bibitem{Paper2}
Washburn, J. (2026). The standard genetic code is globally optimal under
position-weighted mutation cost. \textit{In preparation}.

\bibitem{Paper3}
Washburn, J. (2026). Z-parity and the hierarchy of conservation in the
genetic code. \textit{In preparation}.

\bibitem{Paper4}
Washburn, J. (2026). Toward a zero-junk genome: Recognition Science
predictions for DNA as Z-invariant storage. \textit{In preparation}.

\bibitem{Paper6}
Washburn, J. (2026). Sex as phase-debt cancellation: a formal derivation
of recombination advantage. \textit{In preparation}.

\bibitem{Paper7}
Washburn, J. (2026). Predicting intron placement from helical phase
stress. \textit{In preparation}.

\bibitem{Paper8}
Washburn, J. (2026). The molecular gearbox: a $\varphi$-selective
architecture for biological timing. \textit{In preparation}.

\bibitem{Washburn2025}
Washburn, J. (2025). The Algebra of Reality: A Recognition Science
Derivation of Physical Law. \textit{Axioms} \textbf{15}(2), 90.
\end{thebibliography}

\end{document}
